% ===========================================================================
\section{Introduction}
\label{sec:introduction}
% ===========================================================================

There are two ways to build a road transport inventory and they fail differently.

The bottom-up approach estimates the vehicle fleet, assigns annual kilometres and fuel economy
by class, and sums. It produces a rich result---emissions by vehicle type, useful for
projecting the effect of a fleet policy---and a fuel demand that will not match the national
petroleum supply. The compiler then adjusts something until it does.

The fuel-anchored approach starts from the supply figure and distributes it across classes.
It produces a total that matches national statistics by construction and a composition that is
only as good as the fleet model. Nothing needs reconciling because nothing was independently
estimated.

Which is better depends on which number is more trustworthy. In Pakistan that is not close:
the petroleum supply figure is a metered, taxed, nationally reported quantity, while the
vehicle fleet is a registration database that counts vehicles administratively rather than
by where they operate, includes a residual catch-all class, and says nothing about how far
each is driven. This paper quantifies the gap directly: the bottom-up fuel demand the same
fleet data imply, under declared engineering intensity bands, misses the observed diesel
control by half and overstates metered CNG threefold
(Section~\ref{sec:validation})---errors in both directions, which is the strongest form
of the case for letting fuel set the magnitude.

\subsection{What this paper does}

We report Pakistan's road transport and off-road mobile machinery for FY2023--24 on the
fuel-anchored basis, and treat the consequences of that choice as the subject rather than as
an implementation detail.

The consequence that matters most is that anchoring makes a boundary problem explicit.
Pakistan's petroleum statistics do publish a sectoral split---OCAC's customer-class sales
tables, which this edition anchors on directly---but a customer class is a sales channel,
not an end use: the retail Road row still contains diesel bought at pumps for stationary
agricultural plant and captive generators. This paper reads the observed customer-class
partition into a two-tier fuel-claimant ledger, carves out what a sibling sector measures
(household-generator petrol), partitions off-road machinery, cross-checks the rail
boundary against an independent source, and states plainly what remains inside the road
residual and in which direction it biases the result.

\subsection{Why Pakistan, and why anchoring in particular}

Pakistan's most recent full national inventory submission reports for FY2017--18
\citep{pakbur1}, and the international products that fill the gap must be compared with the
years matched, because Pakistan's fuel demand moved sharply after 2022. On a scope
inferred from EDGAR's category-resolved Pakistani source files---road, rail and civil
aviation---this inventory's corresponding construction is \num{0.85} of EDGAR's
fiscal-year blend of \SI{48.13}{\mega\tonne} \citep{edgar2025booklet}
(Section~\ref{sec:validation}), sitting \emph{under} the benchmark for the reasons the
customer-class boundary predicts. The blend is half of EDGAR's 2023 column and half of
its 2024 one, because this sector's July--June reference period falls across both; the
comparison is a scope-approximate diagnostic rather than a validation, because EDGAR's
transport boundary is inferred rather than published.

Anchoring is particularly appropriate here for a reason specific to the country's fleet.
Two-wheelers are the largest single class in Pakistan's road transport---under the
reference allocation they receive \SI{13.40}{\mega\tonne}, \SI{34.1}{\percent} of
on-road \COtwo{}---and their annual kilometres are among the least well-constrained
inputs in this Pakistani fleet model. A bottom-up estimate would rest a very large term on
one of its weakest inputs. Anchored, that same uncertainty affects only how the fuel is
split between classes, which is exactly what happened when this edition corrected the
split: two-wheelers fell by six megatonnes, heavy diesel rose past them, and the sector
total moved by \SI{0.05}{\mega\tonne} (Section~\ref{sec:results}).

\subsection{Contributions}

\begin{enumerate}[leftmargin=*,itemsep=2pt,topsep=3pt]
\item \textbf{A fuel-anchored road transport inventory for Pakistan} anchored on the
  customer-class rows of the national sectoral sales table, so the sector total
  reconciles with petroleum statistics by construction and the bulk non-road claimants
  are excluded by the source's own structure (Section~\ref{sec:methods}).
\item \textbf{A two-tier fuel-claimant ledger, mostly observed}: the national tier is
  OCAC's own exhaustive customer-class partition; the road tier carries one relayed
  carve-out (generator petrol, closed against the residential module), one partition
  (off-road machinery), an independent rail cross-check between two unrelated sources,
  and the unresolved retail claimants as bounded status rows instead of silence
  (Table~\ref{tab:ledger}).
\item \textbf{A quantified instance of the IPCC parallel-table trap}, in which the plausible
  wrong table leaves \COtwo{} exactly right and moves the category by \SI{8}{\percent} through
  \NtwoO{} alone---invisible to every \COtwo{}-dominated comparison available
  (Section~\ref{sec:factors}).
\item \textbf{A source-traceable road air-pollutant estimate} using EMEP/EEA Tier~1 fuel-based exhaust factors, which the guidebook derives from its hot and cold Tier~3 method for a typical 2010 EU-27 fleet (no Pakistan-specific cold-start treatment is applied), separate petrol evaporation and tyre/brake/road-surface wear. Transferred fleet factors, technology, sulfur compliance and resuspended road dust retain their actual limitations (Section~\ref{sec:air}).
\item \textbf{A provincial petroleum allocation from OCAC's monthly area editions},
  allocated per fuel, with a published diagnostic demonstrating the failure of the
  registration proxy it replaces and an allocation crosswalk stating which streams ride
  which proxy (Section~\ref{sec:results}).
\end{enumerate}

Method shared with the other nine sectors of the national inventory---the module contract,
the accounting closures, the parameter-uncertainty treatment---is documented separately
\citep{ahmad2026national} and not repeated here.
